% --- CORRECCIÓN IMPORTANTE EN LA PRIMERA LÍNEA ---
% Agregamos "xcolor=table" para que funcione \rowcolor sin errores
\documentclass[aspectratio=169, 10pt, xcolor=table]{beamer}

% --- Configuración del Tema (estilo profesional como el ejemplo) ---
\usetheme{Madrid}
\usecolortheme{dolphin}

% --- Paquetes necesarios ---
\usepackage[utf8]{inputenc}
\usepackage[spanish]{babel}
\usepackage{amsmath, amssymb}
\usepackage{booktabs}
\usepackage{graphicx}
\usepackage{listings}
\usepackage{tikz}
\usetikzlibrary{arrows.meta, positioning, shapes.geometric}

% --- Ruta de Imágenes (crea la carpeta "imagenes" y coloca ahí tus archivos) ---
\graphicspath{{./imagenes/}}

% --- Estilo para código Python ---
\definecolor{codegreen}{rgb}{0,0.6,0}
\definecolor{codegray}{rgb}{0.5,0.5,0.5}
\definecolor{codepurple}{rgb}{0.58,0,0.82}
\definecolor{backcolour}{rgb}{0.95,0.95,0.92}
\lstset{
    language=Python,
    backgroundcolor=\color{backcolour},
    commentstyle=\color{codegreen},
    keywordstyle=\color{blue},
    numberstyle=\tiny\color{codegray},
    stringstyle=\color{codepurple},
    basicstyle=\ttfamily\scriptsize,
    breaklines=true,
    frame=single,
    showstringspaces=false
}

% --- Pie de página personalizado ---
\makeatletter

\setbeamertemplate{footline}{
  \leavevmode%
  \hbox{%
  \begin{beamercolorbox}[wd=.333333\paperwidth,ht=2.25ex,dp=1ex,center]{author in head/foot}%
    \usebeamerfont{author in head/foot}\insertshortauthor
  \end{beamercolorbox}%
  \begin{beamercolorbox}[wd=.333333\paperwidth,ht=2.25ex,dp=1ex,center]{title in head/foot}%
    \usebeamerfont{title in head/foot}Sesión 14
  \end{beamercolorbox}%
  \begin{beamercolorbox}[wd=.333333\paperwidth,ht=2.25ex,dp=1ex,right]{date in head/foot}%
    \usebeamerfont{date in head/foot}Machine Learning \hfill \insertframenumber/\inserttotalframenumber \hspace*{0.3cm}
  \end{beamercolorbox}}%
  \vskip0pt%
}

\makeatother

% --- Datos de la portada ---
\title[SESIÓN 14]{\textbf{Sesión 14}}
\subtitle{Despliegue de Modelos (MLOps): \\ De Jupyter Notebook a producción}
\author[Prof. C. Sánchez]{Carlos César Sánchez Coronel}
\institute{\textbf{MACHINE LEARNING}}
\date{}

% Logo opcional (descomentar si tienes la imagen)
% \titlegraphic{\includegraphics[height=1.5cm]{logo.png}}

% --- Eliminar la barra de navegación (opcional) ---
\setbeamertemplate{navigation symbols}{}

% --- Índice al inicio de cada sección ---
\AtBeginSection[]{
  \begin{frame}
    \frametitle{Contenido}
    \tableofcontents[currentsection]
  \end{frame}
}

% ============================================================================
% INICIO DEL DOCUMENTO
% ============================================================================
\begin{document}

% --- Portada ---
\begin{frame}
    \titlepage
\end{frame}

% --- Índice general ---
\begin{frame}{Contenido}
    \tableofcontents
\end{frame}

% ============================================================================
% SECCIÓN 1: INTRODUCCIÓN A MLOps
% ============================================================================
\section{Introducción a MLOps}

\begin{frame}{El último kilómetro del machine learning}
    \begin{itemize}
        \item Construir un modelo en Jupyter es solo el 10\% del trabajo.
        \item El 90\% restante es llevarlo a producción y mantenerlo.
        \item \textbf{MLOps}: disciplina que combina ML, ingeniería de software y DevOps.
        \item Objetivos: disponibilidad, escalabilidad, reproducibilidad, monitoreo continuo.
    \end{itemize}
    % Basado en introducción
    % IMAGEN: Ciclo de vida MLOps
\end{frame}

\begin{frame}{Entrenamiento vs Inferencia}
    \begin{columns}[t]
        \column{0.5\textwidth}
        \textbf{Entrenamiento}
        \begin{itemize}
            \item Offline, por lotes.
            \item Puede ser lento (horas/días).
            \item Grandes volúmenes de datos.
            \item Poco frecuente.
        \end{itemize}
        \column{0.5\textwidth}
        \textbf{Inferencia}
        \begin{itemize}
            \item Online (tiempo real) o batch.
            \item Baja latencia (online).
            \item Alta concurrencia.
            \item Continua.
        \end{itemize}
    \end{columns}
    \vspace{0.3cm}
    La arquitectura debe separar ambos pipelines.
    % Basado en Pregunta 1
\end{frame}

% ============================================================================
% SECCIÓN 2: SERIALIZACIÓN DE MODELOS
% ============================================================================
\section{Serialización de Modelos}

\begin{frame}{Formatos de serialización}
    \begin{columns}[t]
        \column{0.33\textwidth}
        \textbf{Pickle}
        \begin{itemize}
            \item Estándar de Python.
            \item Inseguro, problemas de compatibilidad.
        \end{itemize}
        \column{0.33\textwidth}
        \textbf{Joblib}
        \begin{itemize}
            \item Optimizado para arrays NumPy.
            \item Recomendado para scikit-learn.
        \end{itemize}
        \column{0.33\textwidth}
        \textbf{ONNX}
        \begin{itemize}
            \item Formato abierto e interoperable.
            \item Permite optimizaciones de inferencia.
        \end{itemize}
    \end{columns}
    \vspace{0.3cm}
    \textbf{MLflow}: plataforma para gestionar modelos, con registro y versionado.
    % Basado en Preguntas 2 y 3
\end{frame}

\begin{frame}[fragile]{Ejemplos de serialización}
    \begin{lstlisting}
import pickle
with open('modelo.pkl', 'wb') as f:
    pickle.dump(model, f)

import joblib
joblib.dump(model, 'modelo.joblib')
model = joblib.load('modelo.joblib')

# ONNX (skl2onnx)
from skl2onnx import convert_sklearn
from skl2onnx.common.data_types import FloatTensorType
initial_type = [('float_input', FloatTensorType([None, X.shape[1]]))]
onx = convert_sklearn(model, initial_types=initial_type)
with open("model.onnx", "wb") as f:
    f.write(onx.SerializeToString())
    \end{lstlisting}
\end{frame}

% ============================================================================
% SECCIÓN 3: CREACIÓN DE APIs CON FLASK Y FASTAPI
% ============================================================================
\section{Creación de APIs con Flask y FastAPI}

\begin{frame}{Flask vs FastAPI}
    \begin{table}
        \centering
        \begin{tabular}{l|c|c}
            \toprule
            \textbf{Característica} & \textbf{Flask} & \textbf{FastAPI} \\
            \midrule
            Rendimiento & Bajo & Alto (asíncrono) \\
            Documentación automática & No & Sí (Swagger, ReDoc) \\
            Validación de datos & Con extensiones & Integrada (Pydantic) \\
            Soporte asíncrono & Limitado & Nativo \\
            \bottomrule
        \end{tabular}
    \end{table}
    % Basado en Pregunta 5
\end{frame}

\begin{frame}[fragile]{Ejemplo con Flask}
    \begin{lstlisting}
from flask import Flask, request, jsonify
import joblib

app = Flask(__name__)
model = joblib.load('modelo.joblib')

@app.route('/predict', methods=['POST'])
def predict():
    data = request.get_json()
    features = data['features']
    pred = model.predict([features])
    return jsonify({'prediction': pred.tolist()})

if __name__ == '__main__':
    app.run(host='0.0.0.0', port=5000)
    \end{lstlisting}
\end{frame}

\begin{frame}[fragile]{Ejemplo con FastAPI}
    \begin{lstlisting}
from fastapi import FastAPI
from pydantic import BaseModel
import joblib

app = FastAPI()
model = joblib.load('modelo.joblib')

class Item(BaseModel):
    features: list

@app.post('/predict')
def predict(item: Item):
    pred = model.predict([item.features])
    return {'prediction': pred.tolist()}
    \end{lstlisting}
    \begin{itemize}
        \item Documentación interactiva en \texttt{/docs}.
        \item Validación automática con Pydantic.
    \end{itemize}
    % Basado en Preguntas 5 y 6
\end{frame}

% ============================================================================
% SECCIÓN 4: CONTENERIZACIÓN CON DOCKER
% ============================================================================
\section{Contenerización con Docker}

\begin{frame}{¿Por qué Docker?}
    \begin{itemize}
        \item Elimina el problema de ``funciona en mi máquina''.
        \item Empaqueta la aplicación con todas sus dependencias.
        \item Misma ejecución en cualquier entorno.
    \end{itemize}
    % Basado en Pregunta 7
\end{frame}

\begin{frame}[fragile]{Dockerfile para una API de modelo}
    \begin{lstlisting}
FROM python:3.9-slim

WORKDIR /app

COPY requirements.txt .
RUN pip install -r requirements.txt

COPY model.joblib .
COPY app.py .

CMD ["uvicorn", "app:app", "--host", "0.0.0.0", "--port", "80"]
    \end{lstlisting}
    \begin{itemize}
        \item Construir: \texttt{docker build -t modelo-api .}
        \item Ejecutar: \texttt{docker run -p 80:80 modelo-api}
    \end{itemize}
    % Basado en Pregunta 7
\end{frame}

% ============================================================================
% SECCIÓN 5: VERSIONAMIENTO DE MODELOS Y DATOS
% ============================================================================
\section{Versionamiento de Modelos y Datos}

\begin{frame}{MLflow para seguimiento de experimentos}
    \begin{block}{Componentes principales}
        \begin{itemize}
            \item \textbf{Tracking}: registro de parámetros, métricas, artefactos.
            \item \textbf{Projects}: empaquetado de código.
            \item \textbf{Models}: formato estándar para modelos.
            \item \textbf{Model Registry}: almacén centralizado con versionado y transiciones (staging, production).
        \end{itemize}
    \end{block}
    % Basado en Pregunta 3
\end{frame}

\begin{frame}[fragile]{Ejemplo con MLflow}
    \begin{lstlisting}
import mlflow

with mlflow.start_run():
    mlflow.log_param("n_estimators", 100)
    mlflow.log_metric("auc", 0.85)
    mlflow.sklearn.log_model(model, "model")

# Cargar una versión específica
model = mlflow.pyfunc.load_model("models:/credit_model/1")
    \end{lstlisting}
\end{frame}

\begin{frame}{DVC (Data Version Control)}
    \begin{itemize}
        \item Versiona datasets y modelos (archivos grandes).
        \item Integrado con Git: Git guarda metadatos, DVC maneja el almacenamiento remoto.
        \item Flujo: \texttt{dvc add data.csv} → \texttt{git add data.csv.dvc} → \texttt{git commit} → \texttt{dvc push}
    \end{itemize}
    % Basado en Pregunta 4
\end{frame}

% ============================================================================
% SECCIÓN 6: MONITOREO DE MODELOS EN PRODUCCIÓN
% ============================================================================
\section{Monitoreo de Modelos en Producción}

\begin{frame}{Data drift y concept drift}
    \begin{itemize}
        \item \textbf{Data drift}: cambio en la distribución de las variables de entrada.
        \item \textbf{Concept drift}: cambio en la relación entrada-salida.
        \item Causas: cambios en el mercado, nuevos hábitos, crisis económicas.
    \end{itemize}
    % Basado en Pregunta 8
\end{frame}

\begin{frame}{Population Stability Index (PSI)}
    \[
    PSI = \sum_{i=1}^{B} (p_i - q_i) \cdot \ln\left(\frac{p_i}{q_i}\right)
    \]
    donde $p_i$ = proporción en producción, $q_i$ = proporción en entrenamiento.
    \begin{itemize}
        \item $PSI < 0.1$: cambio pequeño.
        \item $0.1 \le PSI < 0.25$: cambio moderado.
        \item $PSI \ge 0.25$: cambio significativo (reentrenar).
    \end{itemize}
    % Basado en Pregunta 9
\end{frame}

\begin{frame}[fragile]{Cálculo de PSI en Python}
    \begin{lstlisting}
import numpy as np

def psi(expected, actual, bins=10):
    expected = np.clip(expected, 0.001, 0.999)  # evitar log(0)
    actual = np.clip(actual, 0.001, 0.999)
    return np.sum((actual - expected) * np.log(actual / expected))
    \end{lstlisting}
\end{frame}

\begin{frame}{Test de Kolmogorov-Smirnov (KS)}
    \begin{itemize}
        \item Compara dos distribuciones continuas sin discretizar.
        \item Estadístico: máxima diferencia entre las funciones de distribución acumulada.
        \item p-valor bajo $\Rightarrow$ distribuciones diferentes.
        \item Útil para variables continuas.
    \end{itemize}
    % Basado en Pregunta 10
\end{frame}

\begin{frame}{Herramientas de monitoreo}
    \begin{itemize}
        \item \textbf{Evidently AI}: reportes de drift, calidad de datos, rendimiento.
        \item \textbf{WhyLabs / Whylogs}: perfilado de datos y detección de drift.
        \item \textbf{Great Expectations}: validación de datos (expectativas).
        \item \textbf{Prometheus + Grafana}: métricas del sistema.
    \end{itemize}
    % Basado en Pregunta 11
\end{frame}






% ============================================================================
% ESTRATEGIAS DE DESPLIEGUE Y EXPERIMENTACIÓN: A/B TESTING
% ============================================================================
\section{Estrategias de Despliegue y A/B Testing}

\begin{frame}{A/B Testing para modelos en producción}
    \begin{block}{¿Qué es un A/B test?}
        Experimento controlado donde se comparan dos versiones de un modelo (A = control, B = tratamiento) asignando usuarios aleatoriamente.
    \end{block}
    \begin{itemize}
        \item \textbf{Objetivo}: determinar si el nuevo modelo (B) mejora métricas de negocio respecto al actual (A).
        \item \textbf{Asignación aleatoria}: asegura que los grupos sean comparables.
        \item \textbf{Métrica primaria}: una única métrica clave (ej. tasa de conversión, ingresos por usuario).
    \end{itemize}
    % Basado en Pregunta 13 del solucionario
\end{frame}

\begin{frame}{Diseño de un A/B test para modelos}
    \begin{enumerate}
        \item \textbf{Definir la hipótesis}: ej. "el nuevo modelo de recomendación aumenta el tiempo de visualización en un 5\%".
        \item \textbf{Elegir métricas}: primaria (tiempo de visualización) y secundarias (clic, retención).
        \item \textbf{Calcular tamaño de muestra} necesario para significancia estadística (poder, nivel de confianza).
        \item \textbf{Asignar usuarios aleatoriamente} (por ejemplo, mediante hash de usuario).
        \item \textbf{Ejecutar el experimento} durante un período representativo (evitar estacionalidades).
        \item \textbf{Analizar resultados}: test estadístico (t-test, bootstrap) y significancia práctica (no solo p-valor).
    \end{enumerate}
\end{frame}

\begin{frame}{Métricas de negocio vs técnicas}
    \begin{columns}[t]
        \column{0.5\textwidth}
        \textbf{Métricas técnicas}
        \begin{itemize}
            \item Precisión, recall, AUC
            \item RMSE, MAE
            \item Latencia, throughput
        \end{itemize}
        \column{0.5\textwidth}
        \textbf{Métricas de negocio}
        \begin{itemize}
            \item Tasa de conversión
            \item Ingresos por usuario
            \item Retención / churn
            \item Tiempo de sesión
        \end{itemize}
    \end{columns}
    \vspace{0.3cm}
    \begin{alertblock}{Importante}
        Una mejora técnica no siempre se traduce en mejora de negocio. El A/B test valida el impacto real.
    \end{alertblock}
\end{frame}

\begin{frame}{Ejemplo: comparar dos modelos de recomendación}
    \begin{block}{Contexto}
        Plataforma de streaming quiere probar un nuevo algoritmo de recomendación.
    \end{block}
    \begin{itemize}
        \item Grupo A (control): modelo actual (basado en popularidad).
        \item Grupo B (tratamiento): nuevo modelo (filtrado colaborativo).
        \item Duración: 2 semanas, 100,000 usuarios por grupo.
        \item Métrica primaria: tiempo medio de visualización por usuario.
        \item Resultado: B muestra un aumento del 3\% con p-valor < 0.01.
        \item Decisión: desplegar B gradualmente (canary).
    \end{itemize}
\end{frame}

\begin{frame}{Precauciones estadísticas}
    \begin{itemize}
        \item \textbf{Significancia estadística}: usar p-valor (ej. < 0.05) y considerar corrección por múltiples comparaciones si hay varias métricas.
        \item \textbf{Tamaño de muestra}: calcular previamente para evitar falsos negativos (poder estadístico).
        \item \textbf{Efecto mínimo detectable}: definir qué mejora es relevante para el negocio.
        \item \textbf{Duración}: suficiente para capturar variabilidad (evitar sesgos por día de la semana).
        \item \textbf{Interferencia entre grupos}: asegurar que los grupos no interactúen (ej. red social).
    \end{itemize}
\end{frame}

\begin{frame}{Integración con despliegue canario}
    \begin{itemize}
        \item El A/B test puede ser la primera fase de un despliegue canario.
        \item \textbf{Canario}: enrutar un pequeño % de tráfico al nuevo modelo.
        \item \textbf{Medir} las mismas métricas de negocio durante el canario.
        \item Si los resultados son positivos, aumentar gradualmente el porcentaje.
        \item Si no, revertir automáticamente (rollback).
    \end{itemize}
    % Basado en Pregunta 12
\end{frame}





% ============================================================================
% SECCIÓN 7: CASO INTEGRADO - MODELO DE RIESGO CREDITICIO
% ============================================================================
\section{Caso Integrado: Modelo de Riesgo Crediticio}

\begin{frame}{Paso 1: Entrenar y guardar el modelo}
    \begin{lstlisting}
import pandas as pd
from sklearn.ensemble import RandomForestClassifier
import joblib

df = pd.read_csv('credito.csv')
X = df.drop('incumple', axis=1)
y = df['incumple']
model = RandomForestClassifier(n_estimators=100)
model.fit(X, y)
joblib.dump(model, 'credit_model.joblib')
    \end{lstlisting}
    % Basado en teoría
\end{frame}

\begin{frame}[fragile]{Paso 2: Crear API con FastAPI}
    \begin{lstlisting}
from fastapi import FastAPI
from pydantic import BaseModel
import joblib
import numpy as np

app = FastAPI()
model = joblib.load('credit_model.joblib')

class Applicant(BaseModel):
    edad: int
    ingresos: float
    deuda: float
    antiguedad_laboral: int

@app.post('/predict')
def predict(applicant: Applicant):
    features = np.array([[applicant.edad, applicant.ingresos,
                          applicant.deuda, applicant.antiguedad_laboral]])
    proba = model.predict_proba(features)[0, 1]
    return {'probabilidad_incumplimiento': proba}
    \end{lstlisting}
\end{frame}

\begin{frame}{Pasos 3-5: Docker, MLflow y monitoreo}
    \begin{itemize}
        \item \textbf{Docker}: contenerizar la API con el Dockerfile visto.
        \item \textbf{MLflow}: registrar el modelo y versionarlo.
        \item \textbf{Monitoreo}: calcular PSI semanal de las características y predicciones; alertar si PSI > 0.25.
    \end{itemize}
    % Basado en teoría
\end{frame}

% ============================================================================
% SECCIÓN 8: REVISIÓN GENERAL DEL CURSO
% ============================================================================
\section{Revisión General del Curso}

\begin{frame}{Mapa conceptual del curso}
    % IMAGEN: Diagrama de integración de todas las sesiones
    \begin{center}
        \begin{tikzpicture}[node distance=0.8cm, font=\tiny]
            \node[draw, fill=blue!20, rectangle] (s1) {S1-2: Matemáticas};
            \node[draw, fill=blue!20, below of=s1] (s3) {S3: EDA y Feature Eng};
            \node[draw, fill=green!20, right of=s3, xshift=1.5cm] (s4) {S4: Regresión Lineal};
            \node[draw, fill=green!20, below of=s4] (s5) {S5: Clasif I (Logística)};
            \node[draw, fill=green!20, below of=s5] (s6) {S6: Clasif II (KNN, NB, SVM)};
            \node[draw, fill=orange!20, right of=s5, xshift=1.5cm] (s7) {S7: Árboles y RF};
            \node[draw, fill=orange!20, below of=s7] (s8) {S8: Gradient Boosting};
            \node[draw, fill=red!20, right of=s8, xshift=1.5cm] (s9) {S9: Validación};
            \node[draw, fill=red!20, below of=s9] (s10) {S10: No Supervisado};
            \node[draw, fill=red!20, below of=s10] (s11) {S11: Series Temporales};
            \node[draw, fill=purple!20, right of=s10, xshift=1.5cm] (s12) {S12: Complementarios};
            \node[draw, fill=purple!20, below of=s12] (s13) {S13: Recomendación};
            \node[draw, fill=purple!20, below of=s13] (s14) {S14: Interpretabilidad};
            \node[draw, fill=brown!20, right of=s13, xshift=1.5cm] (s15) {S15: MLOps};
            \draw[->] (s1) -- (s3);
            \draw[->] (s3) -- (s4);
            \draw[->] (s4) -- (s5);
            \draw[->] (s5) -- (s6);
            \draw[->] (s5) -- (s7);
            \draw[->] (s7) -- (s8);
            \draw[->] (s6) -- (s9);
            \draw[->] (s8) -- (s9);
            \draw[->] (s9) -- (s10);
            \draw[->] (s10) -- (s11);
            \draw[->] (s10) -- (s12);
            \draw[->] (s12) -- (s13);
            \draw[->] (s13) -- (s14);
            \draw[->] (s14) -- (s15);
        \end{tikzpicture}
    \end{center}
\end{frame}

% ============================================================================
% SECCIÓN 9: CONEXIÓN CON EL FUTURO PROFESIONAL
% ============================================================================
\section{Conexión con el Futuro Profesional}

\begin{frame}{Más allá del curso}
    \begin{itemize}
        \item Dominar MLOps diferencia a un junior de un senior.
        \item Habilidades valoradas: despliegue, monitoreo, CI/CD, escalabilidad.
        \item Herramientas a explorar: Kubeflow, Airflow, MLflow, DVC, Kubernetes.
        \item Práctica continua, lectura de papers, competencias (Kaggle).
        \item Mantenerse actualizado con nuevas técnicas.
    \end{itemize}
\end{frame}

% ============================================================================
% SECCIÓN 10: RESUMEN Y CONCLUSIÓN
% ============================================================================
\section{Resumen y Conclusión}

\begin{frame}{Lo que hemos aprendido}
    \begin{exampleblock}{Conceptos clave}
        \begin{itemize}
            \item MLOps: integración de ML con operaciones.
            \item Serialización: pickle, joblib, ONNX, MLflow.
            \item APIs: Flask, FastAPI (recomendado).
            \item Contenerización: Docker.
            \item Versionado: MLflow (modelos), DVC (datos).
            \item Monitoreo: PSI, KS, herramientas (Evidently, Prometheus).
            \item Estrategias de despliegue: canary, blue-green.
            \item A/B testing y seguridad.
        \end{itemize}
    \end{exampleblock}
\end{frame}

\begin{frame}{Conclusión del curso}
    \begin{itemize}
        \item Recorrido completo: de matemáticas a producción.
        \item Modelos supervisados, no supervisados, series temporales, recomendación.
        \item Interpretabilidad y MLOps como cierre.
        \item El aprendizaje continúa: práctica, actualización, comunidad.
    \end{itemize}
    \vspace{1cm}
    \centering
    \Large{¡Enhorabuena por llegar hasta el final!}
\end{frame}

% --- Diapositiva final ---
\begin{frame}
    \centering
    \Huge \textbf{¡Gracias!}

\end{frame}

\end{document}